ModeBench makes a missing part of evaluation measurable: the range of distinct,
verified outcomes a model can produce for one problem. Accuracy tells us whether
a model can succeed; diversity helps describe the alternatives still available
to it. That distinction matters because different approaches can have different
strengths. One may succeed where another fails, or remain useful when the
problem or its constraints change. Knowing several ways to solve could therefore
be valuable even when a familiar approach already works.

Our experiments show that replay complements learning from fresh samples.
Remembering verified solutions improves held-out correctness and preserves
more correct output alternatives across the completed comparisons. ModeBench
lets us measure those alternatives through executable, task-defined identities.
It provides a controlled starting point for understanding what training
preserves, alongside what training learns to solve.

\textbf{Limits and next steps.} ModeBench provides a controlled setting for
studying diverse solutions. Understanding how this diversity relates to broader
mathematical reasoning and later problem solving remains an open direction. Replay can
retain only modes it discovers and stores. Individual-mode survival and
size-dependent update geometry remain unmeasured
(Appendices~\ref{app:ablations} and~\ref{app:theory-size}); the optimizer and
admission results are conditional. Theory studies idealized mean flows.
We have not shown that broader support
causes the accuracy gains. The MaxRL factorial has eleven complete paired
blocks; Falcon Countdown's Dr.GRPO replay comparison has four admissible seeds,
and only Graph has a complete Level-2 factorial. Richer mathematical tasks
should test which alternatives survive training, when a different approach
succeeds after another fails, and whether preserving those options improves
problem solving.

